




% low code

\cite{alsarraj2026bridging} - 
a key challenge in the industry is the translation of high level models to executable code. They present a study that introduces a visual tool that auto converts class domain models to executable code. Mentions the benefits of developing this is that its a major software engineerring advancement, automtes the development of a high level design concept into executabkole code that is accurate and consistent with design specification . Has consistency, maintainability, and scalability throughout the software lifecycle


% privacy 


% engineering implications
\cite{cruz2009assessing}
an engineering model must possess the following five key
quality characteristics to a sufficient degree \cite{selic2003pragmatics} abstraction, understandability,
accuracy, predictiveness and inexpensiveness. They specifically focus on assessing the understandability of statecharts for practitioners, which is tied to dependability's maintainability
They specifically looked into composite states
They carried out 3 experiements, replicated the first two, with students from different universities specifically computer science undergrad and graduate
Found the use of compososite sattes actually had a negative effect on the understandability of the diagrams, specifcally for those with a lack of experuence in modelling. However they found when giving complicated task it was much easier to perform/useful to use composite states



% future work - extensions to modelling
\cite{ramos2015extending}
extends state charts to model interactions in system of systems. Propose a notartion to represent othogonal states deriving from \cite{harel1987statecharts}, and to represent systems states similar in strcuture but belonging to different systems, and then bc their orthogonal they can model communication among parallel states to represent system interactions. Harel actually attepted to show an idea where multiple othoginal states withj the same structure can intreract to achieve a specific goal. \cite{zave2002component} attempted this problem too in relation to modelling telecommunications systems. They noted its impractical to model all the detail in a single state chart as it is now.
Could be mentioned as a future point, currently we use state charts to model contributions being allowed into the system but what about also extending so these state charts can be mapped together to represent communication
For them they represent every constituent using a generic state, and then instances correspond to parameterized states
Presented a symbolic notation to make the above by extending state chart semantics